
\chapter{$\mathcal{W}_1$ 最適輸送}\label{chap:w1}

本章では，地上コストが距離そのもの $c(x,y) = d(x,y)$ である場合，
すなわち $p = 1$ の Wasserstein 距離 $\mathcal{W}_1$ を扱う．
$\mathcal{W}_1$ は Monge が 1781 年に提起した
原初的な設定に対応する距離であり，
コンピュータビジョンでは
\textbf{Earth Mover's Distance}~\cite{Rubner2000}
の名で広く知られている．

$\mathcal{W}_2$（2乗コスト）と比較して，
$\mathcal{W}_1$ には以下の特徴がある：
\begin{itemize}
  \item 双対問題が \textbf{Lipschitz 制約}に帰着され，
    簡潔で美しい形を持つ．
  \item 符号付き測度（正とは限らない）に対しても定義でき，
    \textbf{ノルム}を定める．
  \item 最適カップリングの一意性が一般には成り立たないため，
    Monge 写像の構造理論は $\mathcal{W}_2$ より複雑になる．
  \item グラフ上では\textbf{最小費用流問題}に帰着され，
    効率的なアルゴリズムが利用できる．
\end{itemize}


% ============================================================
\section{距離空間上の $\mathcal{W}_1$}\label{sec:w1-metric}

\subsection{Kantorovich--Rubinstein 双対性}\label{subsec:kr-duality}

$\mathcal{X} = \mathcal{Y}$ を距離空間，
$d$ をその距離関数とする．
$c(x, y) = d(x, y)$ のとき，
Kantorovich 双対問題（式~\eqref{eq:dual-general}）は
$c$変換の性質により顕著に簡略化される．

関数 $f \in \mathcal{C}(\mathcal{X})$ の
\textbf{Lipschitz 定数}を
\[
  \mathrm{Lip}(f) \defeq
  \sup_{\substack{x, y \in \mathcal{X} \\ x \neq y}}
  \frac{|f(x) - f(y)|}{d(x, y)}
\]
で定義する．

\begin{proposition}{$c$変換と Lipschitz 条件}{c-transform-lip}
$c(x, y) = d(x, y)$ のとき，
$f = g^c$ ならば $\mathrm{Lip}(f) \leq 1$ であり，
逆に $\mathrm{Lip}(f) \leq 1$ ならば $f^c = -f$ が成り立つ．
\end{proposition}

\begin{proof}
$f = g^c$ とすると，
$f(x) = \inf_{z} d(x, z) - g(z)$ であるから，
任意の $x, y$ に対し
\begin{align*}
  |f(x) - f(y)|
  &= \left|\inf_z d(x,z) - g(z) - \inf_z d(y,z) + g(z)\right| \\
  &\leq \sup_z |d(x,z) - d(y,z)| \leq d(x,y)
\end{align*}
より $\mathrm{Lip}(f) \leq 1$．

逆に $\mathrm{Lip}(f) \leq 1$ のとき，$g = -f$ とおくと
$f(y) - d(x,y) \leq f(x) \leq f(y) + d(x,y)$ から
$f^c(y) = \inf_x d(x,y) - f(x) = -f(y)$ が確認できる．
\end{proof}

この命題から，双対問題を単一変数 $f$ で書くと
次の\textbf{Kantorovich--Rubinstein 公式}が得られる：

\begin{theorem}{Kantorovich--Rubinstein 双対性}{kr-duality}
\begin{equation}\label{eq:kr-dual}
  \mathcal{W}_1(\alpha, \beta)
  = \max_{f}\; \left\{
    \int_{\mathcal{X}} f(x)\, (\mathrm{d}\alpha(x) - \mathrm{d}\beta(x))
    \;:\; \mathrm{Lip}(f) \leq 1
  \right\}.
\end{equation}
\end{theorem}

この公式は，$\mathcal{W}_1$ が
\textbf{1-Lipschitz 関数全体を判別関数とするときの}
$\alpha$ と $\beta$ の最大期待値差であることを示している．

\subsection{$\mathcal{W}_1$ はノルムを定める}\label{subsec:w1-norm}

式~\eqref{eq:kr-dual}の右辺は $\alpha - \beta$
（符号付き測度）の線形汎関数の上限であるから，
$\mathcal{W}_1(\alpha, \beta) = \|\alpha - \beta\|_{\mathcal{W}_1}$
と書ける．すなわち $\mathcal{W}_1$ は（全質量が等しい）
符号付き測度の空間上の\textbf{ノルム}を定める．
この性質は $p > 1$ の $\mathcal{W}_p$ にはなく，
$\mathcal{W}_1$ に特有の構造である．

\subsection{離散的な定式化}\label{subsec:w1-discrete}

離散測度
$\alpha - \beta = \sum_k \mathbf{m}_k \delta_{z_k}$
（$\sum_k \mathbf{m}_k = 0$）に対し，
式~\eqref{eq:kr-dual}は有限次元の最適化問題になる：
\begin{equation}\label{eq:w1-discrete}
  \mathcal{W}_1(\alpha, \beta)
  = \max_{(\mathbf{f}_k)_k} \left\{
    \sum_k \mathbf{f}_k \mathbf{m}_k
    \;:\;
    \forall\, (k, \ell),\;\;
    |\mathbf{f}_k - \mathbf{f}_\ell| \leq d(z_k, z_\ell)
  \right\}.
\end{equation}
これは二次錐制約（$|\cdot| \leq \cdot$）を持つ凸計画であり，
内点法や近接法で解ける．


% ============================================================
\section{ユークリッド空間上の $\mathcal{W}_1$}\label{sec:w1-euclidean}

$\mathcal{X} = \mathcal{Y} = \R^d$，$d(x,y) = \|x - y\|$ のとき，
Lipschitz 制約はポテンシャル $f$ の\textbf{勾配ノルム}の制約に局所化できる：

\begin{equation}\label{eq:w1-gradient}
  \mathcal{W}_1(\alpha, \beta)
  = \max_f\; \left\{
    \int_{\R^d} f(x)\, (\mathrm{d}\alpha - \mathrm{d}\beta)(x)
    \;:\; \|\nabla f\|_\infty \leq 1
  \right\}.
\end{equation}

この双対問題に対応する主問題は，
\textbf{Beckmann の定式化}と呼ばれる
ベクトル場 $s : \R^d \to \R^d$ に関する最適化問題である：

\begin{equation}\label{eq:beckmann}
  \mathcal{W}_1(\alpha, \beta)
  = \min_s\; \left\{
    \int_{\R^d} \|s(x)\|_2\, \mathrm{d}x
    \;:\; \mathrm{div}(s) = \alpha - \beta
  \right\}.
\end{equation}

ここでベクトル場 $s(x)$ は質量の局所的な移動方向と速度を表し，
発散条件 $\mathrm{div}(s) = \alpha - \beta$ は質量の保存を保証する．
$\alpha$ の台の外で $\mathrm{div}(s) = 0$ であるから，
$s$ はソース（$\alpha$）からシンク（$\beta$）へ質量を流す
\textbf{流れ場}と解釈できる．

Beckmann の定式化は非滑らか凸最適化問題であり，
有限要素法で離散化した後，
Douglas--Rachford 分裂法や ADMM 等の一階法で解かれる．


% ============================================================
\section{グラフ上の $\mathcal{W}_1$}\label{sec:w1-graph}

$\mathcal{X}$ がグラフ $(\{1, \ldots, n\}, \mathcal{E})$ で
辺重み $\mathbf{w}_{i,j} > 0$ を持つとき，
ユークリッド空間の勾配・発散の概念をグラフ上に離散化することで，
$\mathcal{W}_1$ を最小費用流問題として定式化できる．

\subsection{グラフ上の勾配と発散}\label{subsec:graph-operators}

グラフの\textbf{勾配作用素} $\nabla : \R^n \to \R^{\mathcal{E}}$ を
\[
  \forall\, (i, j) \in \mathcal{E}, \quad
  (\nabla \mathbf{f})_{i,j} \defeq \mathbf{f}_i - \mathbf{f}_j
\]
で定義する．その随伴作用素（\textbf{発散}）
$\mathrm{div} : \R^{\mathcal{E}} \to \R^n$ は
\[
  \forall\, i, \quad
  \mathrm{div}(\mathbf{s})_i \defeq
  \sum_{j:\, (i,j) \in \mathcal{E}} (\mathbf{s}_{i,j} - \mathbf{s}_{j,i})
\]
で定義される．

\subsection{主問題と双対問題}\label{subsec:w1-graph-formulation}

式~\eqref{eq:w1-gradient}と~\eqref{eq:beckmann}の
グラフ上での対応は以下のようになる．

\paragraph{双対問題（Lipschitz 制約）.}
\begin{equation}\label{eq:w1-graph-dual}
  \mathrm{W}_1(\mathbf{a}, \mathbf{b})
  = \max_{\mathbf{f} \in \R^n}\; \left\{
    \sum_{i=1}^n \mathbf{f}_i (\mathbf{a}_i - \mathbf{b}_i)
    \;:\; \forall\, (i,j) \in \mathcal{E},\;\;
    |(\nabla \mathbf{f})_{i,j}| \leq \mathbf{w}_{i,j}
  \right\}.
\end{equation}

\paragraph{主問題（最小費用流）.}
\begin{equation}\label{eq:w1-graph-primal}
  \mathrm{W}_1(\mathbf{a}, \mathbf{b})
  = \min_{\mathbf{s} \in \R^{\mathcal{E}}_+}\; \left\{
    \sum_{(i,j) \in \mathcal{E}}
    \mathbf{w}_{i,j}\, \mathbf{s}_{i,j}
    \;:\; \mathrm{div}(\mathbf{s}) = \mathbf{a} - \mathbf{b}
  \right\}.
\end{equation}

式~\eqref{eq:w1-graph-primal}は
\textbf{最小費用流問題} (minimum-cost flow) そのものであり，
辺 $(i,j)$ を通る流量 $\mathbf{s}_{i,j}$ の
重み付き総量を最小化する．
この問題には専用の高効率ソルバーが存在し，
一般の LP ソルバーを使う Kantorovich 問題より高速に解ける．

\begin{remark}{$n \times n$ のカップリング行列が不要}{no-coupling}
グラフ上の $\mathcal{W}_1$ の計算では，
$n \times n$ のカップリング行列 $\mathbf{P}$ を
陽に構成する必要がない．
辺の数 $|\mathcal{E}|$ に比例するサイズのフロー $\mathbf{s}$ を
最適化すればよく，
スパースなグラフ（$|\mathcal{E}| \ll n^2$）では
大幅な計算量の削減につながる．
\end{remark}

\begin{example}{パスグラフ上の $\mathrm{W}_1$}{w1-path}
$\mathcal{X} = \{1, 2, 3, 4\}$ が
$1$---$2$---$3$---$4$ のパスグラフ（辺重み全て $1$）である場合を考える．
$\mathbf{a} = (1, 0, 0, 0)$，$\mathbf{b} = (0, 0, 0, 1)$ とすると，
質量 $1$ をノード $1$ からノード $4$ へ送る必要がある．
最小費用流は $\mathbf{s}_{1,2} = \mathbf{s}_{2,3} = \mathbf{s}_{3,4} = 1$
（3辺を通過）であり，$\mathrm{W}_1 = 3$．
これはグラフ上の最短経路距離 $d(1, 4) = 3$ に一致する．
\end{example}
